\documentclass[11pt,a4paper]{article}
\usepackage[utf8]{inputenc}
\usepackage[T1]{fontenc}
\usepackage[spanish]{babel}
\usepackage{listings}
\usepackage{booktabs}
\usepackage{enumitem}
\usepackage{hyperref}
\usepackage{geometry}
\geometry{margin=2.5cm}

\lstset{
  language=CSharp,
  basicstyle=\ttfamily\small,
  breaklines=true,
  breakatwhitespace=true,
  showstringspaces=false,
  frame=single,
  columns=fullflexible,
  keepspaces=true
}

\begin{document}

% --- PORTADA ---
\begin{titlepage}
    \centering
    \vspace*{1.5cm}
    
    {\Large \textbf{Clase Práctica 4 de Ingeniería de Software}}\\[2cm]
    
    {\Huge \textbf{Mani Molido 2, ahora más molido}}\\[3cm]
    
    {\large \textbf{Integrantes}}\\[0.5cm]
    \begin{tabular}{cl}
        \toprule
        \# & Nombre completo \\
        \midrule
        1 & Adrián Xavier Brioso Jurado \\
        2 & Alejandro García Ruiz \\
        3 & Ana Laura Hernández Cuétino \\
        4 & Juan Carlos Yern Espinosa \\
        5 & Laura Alonso Rivero \\
        \bottomrule
    \end{tabular}\\[2cm]
    
    {\large \textbf{Contribución individual por ejercicio}}\\[0.5cm]
    \begin{tabular}{ll}
        \toprule
        Integrante & Ejercicio(s) en los que trabajó \\
        \midrule
        Adrián Xavier Brioso Jurado & 6a, 6b \\
        Alejandro García Ruiz & 6c, 6d \\
        Ana Laura Hernández Cuétino & 5a, 5c \\
        Juan Carlos Yern Espinosa & 5b \\
        Laura Alonso Rivero & 5d \\
        \bottomrule
    \end{tabular}\\[2cm]
    
    {\large \textbf{Repositorio público del proyecto}}\\[0.5cm]
    \url{https://github.com/Jcyern/Management-for-the-automatic-generation-of-delivery-routes/tree/Entregas-IS}
    
    \vfill
    
    23 de octubre de 2026
\end{titlepage}

% --- CONTENIDO DEL DOCUMENTO ---

\section*{Ejercicio 5 — Sistema de Gestión de Archivos Pesados}

\subsection*{a)}

El patrón que mejor encaja es el Proxy, más precisamente su variante Virtual Proxy, es un sustituto que implemente la misma interfaz que el objeto real y crea/carga el objeto pesado solo en el primer acceso. El cliente sigue trabajando contra \texttt{IArchivoPesado} y no sabe si detrás hay un archivo ya cargado o un proxy que lo cargará bajo demanda.

\textbf{Clases propuestas:}
\begin{itemize}
    \item \texttt{IArchivoPesado}: interfaz común con operaciones como \texttt{Visualizar()}, \texttt{Procesar()}, \texttt{Nombre}, \texttt{TamanoBytes}.
    \item \texttt{VideoHD}, \texttt{DatasetCientifico}, \texttt{ImagenSatelital}: objetos reales pesados. Su constructor o su método de carga reserva/lee el contenido.
    \item \texttt{ProxyArchivoPesado}: implementa \texttt{IArchivoPesado}, mantiene una referencia perezosa al objeto real y un \texttt{Func<IArchivoPesado>} para crearlo.
    \item \texttt{MetricasProxy}: registra accesos, cantidad de cargas y tiempo de carga sin cambiar la interfaz.
    \item Cliente: usa \texttt{IArchivoPesado}; no conoce las clases concretas ni cuándo se cargan.
\end{itemize}

Una lazy initialization simple dentro del objeto real o del cliente no evita necesariamente crear el objeto real, obliga al cliente a conocer la clase concreta o a repetir lógica de carga y tampoco centraliza métricas ni control de concurrencia. El Virtual Proxy es transparente: el cliente ve la misma interfaz y el proxy decide cuándo instanciar el objeto pesado. Un Protection Proxy controla permisos o derechos de acceso; no está pensado para diferir la creación ni para ahorrar memoria. Podría combinarse con el Virtual Proxy si además se necesitara seguridad, pero no resuelve el problema principal.

\subsection*{b)}

\begin{lstlisting}
using System;
using System.Diagnostics;
using System.Threading;

public interface IArchivoPesado
{
    string Nombre { get; }
    long TamanoBytes { get; }
    void Visualizar();
    void Procesar();
}

public sealed class MetricasProxy
{
    private long _accesos;
    private long _cargas;
    private long _msCargaTotal;

    public void RegistrarAcceso() => Interlocked.Increment(ref _accesos);

    public void RegistrarCarga(long ms)
    {
        Interlocked.Increment(ref _cargas);
        Interlocked.Add(ref _msCargaTotal, ms);
    }

    public long Accesos => Interlocked.Read(ref _accesos);
    public long Cargas => Interlocked.Read(ref _cargas);
    public long MsCargaTotal => Interlocked.Read(ref _msCargaTotal);

    public override string ToString() =>
        $"Accesos={Accesos}, Cargas={Cargas}, TiempoCargaTotal={MsCargaTotal} ms";
}

public sealed class ProxyArchivoPesado : IArchivoPesado
{
    private readonly Func<IArchivoPesado> _cargador;
    private readonly object _candado = new();
    private volatile IArchivoPesado? _real;
    private readonly MetricasProxy _metricas = new();

    public string Nombre { get; }
    public long TamanoBytes { get; }

    public ProxyArchivoPesado(string nombre, long tamanoBytes, Func<IArchivoPesado> cargador)
    {
        Nombre = nombre;
        TamanoBytes = tamanoBytes;
        _cargador = cargador;
    }

    public MetricasProxy Metricas => _metricas;

    private IArchivoPesado ObtenerReal()
    {
        var real = _real;
        if (real != null)
        {
            _metricas.RegistrarAcceso();
            return real;
        }

        lock (_candado)
        {
            if (_real == null)
            {
                var sw = Stopwatch.StartNew();
                _real = _cargador();
                sw.Stop();
                _metricas.RegistrarCarga(sw.ElapsedMilliseconds);
            }

            _metricas.RegistrarAcceso();
            return _real;
        }
    }

    public void Visualizar() => ObtenerReal().Visualizar();
    public void Procesar() => ObtenerReal().Procesar();
}

public sealed class VideoHD : IArchivoPesado
{
    private byte[] _contenido;

    public string Nombre { get; }
    public long TamanoBytes { get; }

    public VideoHD(string nombre, int tamanoBytes)
    {
        Nombre = nombre;
        TamanoBytes = tamanoBytes;
        Thread.Sleep(300);
        _contenido = new byte[tamanoBytes];
    }

    public void Visualizar() => Console.WriteLine($"Reproduciendo video HD: {Nombre}");
    public void Procesar() => Console.WriteLine($"Procesando video HD: {Nombre}");
}

public sealed class DatasetCientifico : IArchivoPesado
{
    private byte[] _contenido;

    public string Nombre { get; }
    public long TamanoBytes { get; }

    public DatasetCientifico(string nombre, int tamanoBytes)
    {
        Nombre = nombre;
        TamanoBytes = tamanoBytes;
        Thread.Sleep(400);
        _contenido = new byte[tamanoBytes];
    }

    public void Visualizar() => Console.WriteLine($"Mostrando dataset: {Nombre}");
    public void Procesar() => Console.WriteLine($"Analizando dataset: {Nombre}");
}

public sealed class ImagenSatelital : IArchivoPesado
{
    private byte[] _contenido;

    public string Nombre { get; }
    public long TamanoBytes { get; }

    public ImagenSatelital(string nombre, int tamanoBytes)
    {
        Nombre = nombre;
        TamanoBytes = tamanoBytes;
        Thread.Sleep(250);
        _contenido = new byte[tamanoBytes];
    }

    public void Visualizar() => Console.WriteLine($"Visualizando imagen satelital: {Nombre}");
    public void Procesar() => Console.WriteLine($"Procesando imagen satelital: {Nombre}");
}
\end{lstlisting}

Las métricas se registran dentro del proxy; \texttt{IArchivoPesado} no cambia.

\subsection*{c)}

Si dos hilos llaman \texttt{Visualizar()} sobre el mismo proxy y el objeto real aún no está cargado, ambos podrían intentar crearlo. Para evitarlo se usa \textbf{doble comprobación con bloqueo}:
\begin{enumerate}
    \item Primer chequeo fuera del \texttt{lock}: si \texttt{\_real} ya existe, se retorna rápido.
    \item Si no existe, se entra al \texttt{lock}.
    \item Segundo chequeo dentro del \texttt{lock}: solo el primer hilo carga; el segundo espera y luego encuentra \texttt{\_real} ya asignado.
    \item \texttt{\_real} se declara \texttt{volatile} para que la asignación sea visible entre hilos.
\end{enumerate}

En Java el equivalente sería \texttt{volatile} + bloque \texttt{synchronized}; en C\# también puede usarse \texttt{Lazy<T>} con \texttt{LazyThreadSafetyMode.ExecutionAndPublication}. Las métricas usan \texttt{Interlocked} para incrementos atómicos. Así se evita carga duplicada y se mantiene la interfaz sin cambios.

\subsection*{d)}

Sin proxy, si se crean los 100 archivos de 1 GB y todos cargan su contenido:

\begin{table}[h]
\centering
\begin{tabular}{lr}
\toprule
Escenario & Memoria aproximada del contenido \\
\midrule
Sin Proxy: 100 archivos cargados & 100 GB \\
Con Proxy: 0 archivos accedidos & \(\sim\)0 GB + metadatos \\
Con Proxy: \(k\) archivos accedidos & \(k\) GB + metadatos \\
Con Proxy: 100 archivos accedidos & 100 GB + sobrecoste del proxy \\
\bottomrule
\end{tabular}
\end{table}

El ahorro real es:
\[
\text{Memoria sin Proxy} = N \times T
\]
\[
\text{Memoria con Proxy} = M \times T + \text{overhead}
\]
donde \(N\) es total de archivos, \(T\) es tamaño por archivo y \(M\) es cantidad de archivos efectivamente cargados. Si solo se usan 10 de 100 archivos de 1 GB, se pasa de \(\sim\)100 GB a \(\sim\)10 GB.

\textbf{Trade-off:} el proxy reduce memoria y tiempo de arranque, pero el primer acceso paga la latencia de carga. Los accesos posteriores son rápidos. Si todos los archivos se acaban usando, el proxy no ahorra memoria y añade una indirección mínima. Las métricas permiten decidir si conviene cargar, cachear o descargar.

\section*{Ejercicio 6 — Sistema de Control de Tráfico Aéreo}

\subsection*{a)}

El patrón que mejor encaja es Mediator. Los aviones no deben conocerse ni comunicarse directamente entre sí, ya que eso crearía una red de dependencias caótica; por el contrario, un Mediator centraliza la coordinación de modo que los aviones envían solicitudes a la torre y la torre decide, ordena y responde, mientras cada avión solo conoce la interfaz del mediador.

\textbf{Clases propuestas:}
\begin{itemize}
    \item \texttt{IMediadorTrafico}: interfaz del mediador con métodos \texttt{Registrar(Avion)} y \texttt{Enviar(Mensaje)}.
    \item \texttt{TorreControl}: mediador concreto que mantiene aviones registrados, colas FIFO por avión y procesa mensajes.
    \item \texttt{Avion}: colega con \texttt{Id}, referencia a \texttt{IMediadorTrafico} y estado.
    \item \texttt{Mensaje}: contiene el tipo (\texttt{Aterrizaje}, \texttt{Despegue}, \texttt{CambioAltitud}) y los datos asociados.
    \item Se pueden crear 3 instancias de \texttt{TorreControl} y asignar cada avión a una torre.
\end{itemize}

Por un lado, Observer sirve para notificar cambios, pero no para coordinar decisiones, resolver conflictos ni garantizar orden FIFO por avión; por otro lado, un Pub/Sub directo entre aviones seguiría dejando la lógica de coordinación distribuida y podría generar caos. En cambio, el Mediator centraliza el protocolo, el estado y el ordenamiento.

\subsection*{b)}

\begin{lstlisting}
using System;
using System.Collections.Concurrent;
using System.Collections.Generic;
using System.Linq;

public enum TipoMensaje
{
    Aterrizaje,
    Despegue,
    CambioAltitud
}

public sealed class Mensaje
{
    public string AvionId { get; }
    public TipoMensaje Tipo { get; }
    public int Altitud { get; }

    public Mensaje(string avionId, TipoMensaje tipo, int altitud = 0)
    {
        AvionId = avionId;
        Tipo = tipo;
        Altitud = altitud;
    }
}

public interface IMediadorTrafico
{
    void Registrar(Avion avion);
    void Enviar(Mensaje mensaje);
}

public sealed class TorreControl : IMediadorTrafico
{
    public string Nombre { get; }

    private readonly ConcurrentDictionary<string, Avion> _aviones = new();
    private readonly ConcurrentDictionary<string, ConcurrentQueue<Mensaje>> _colas = new();
    private readonly ConcurrentDictionary<string, object> _candadosPorAvion = new();

    public TorreControl(string nombre) => Nombre = nombre;

    public void Registrar(Avion avion)
    {
        _aviones[avion.Id] = avion;
        _colas[avion.Id] = new ConcurrentQueue<Mensaje>();
        _candadosPorAvion[avion.Id] = new object();
    }

    public void Enviar(Mensaje mensaje)
    {
        var candado = _candadosPorAvion.GetOrAdd(mensaje.AvionId, _ => new object());

        lock (candado)
        {
            _colas[mensaje.AvionId].Enqueue(mensaje);
            ProcesarCola(mensaje.AvionId);
        }
    }

    private void ProcesarCola(string avionId)
    {
        while (_colas[avionId].TryDequeue(out var mensaje))
        {
            switch (mensaje.Tipo)
            {
                case TipoMensaje.Aterrizaje:
                    Console.WriteLine($"{Nombre}: autoriza aterrizaje de {avionId}");
                    break;
                case TipoMensaje.Despegue:
                    Console.WriteLine($"{Nombre}: autoriza despegue de {avionId}");
                    break;
                case TipoMensaje.CambioAltitud:
                    Console.WriteLine($"{Nombre}: {avionId} cambia altitud a {mensaje.Altitud} pies");
                    break;
            }
        }
    }
}

public sealed class Avion
{
    public string Id { get; }
    private readonly IMediadorTrafico _torre;

    public Avion(string id, IMediadorTrafico torre)
    {
        Id = id;
        _torre = torre;
        _torre.Registrar(this);
    }

    public void SolicitarAterrizaje() => _torre.Enviar(new Mensaje(Id, TipoMensaje.Aterrizaje));
    public void SolicitarDespegue() => _torre.Enviar(new Mensaje(Id, TipoMensaje.Despegue));
    public void CambiarAltitud(int pies) => _torre.Enviar(new Mensaje(Id, TipoMensaje.CambioAltitud, pies));
}
\end{lstlisting}

Creación de 3 torres y 10 aviones:

\begin{lstlisting}
var torres = new List<TorreControl>
{
    new TorreControl("Torre Norte"),
    new TorreControl("Torre Sur"),
    new TorreControl("Torre Este")
};

var aviones = Enumerable.Range(1, 10)
    .Select(i => new Avion($"AV{i:00}", torres[(i - 1) % torres.Count]))
    .ToList();

aviones[0].SolicitarAterrizaje();
aviones[1].CambiarAltitud(12000);
aviones[2].SolicitarDespegue();
\end{lstlisting}

La concurrencia se maneja con un candado por avión. Dos mensajes del mismo avión no se procesan a la vez y mantienen FIFO; mensajes de aviones distintos pueden procesarse concurrentemente. \texttt{ConcurrentQueue} evita condiciones de carrera en la cola.

\subsection*{c)}

Si una torre falla, los aviones que tiene asignados pierden su mediador y, si además es la única torre, se convierte en un punto único de fallo; para añadir redundancia sin romper el patrón conviene que los aviones dependan de \texttt{IMediadorTrafico} y no de \texttt{TorreControl}, de modo que se puedan crear torres de respaldo que implementen la misma interfaz. Una estrategia posible es tener una torre primaria y otra secundaria con réplica de colas y estado, heartbeat y reasignación de aviones a la secundaria cuando la primaria no responda, aunque también puede haber tres torres activas por sectores con respaldo entre ellas; así no se rompe el Mediator porque el avión sigue hablando con un mediador abstracto y solo cambia la instancia concreta. Para no perder mensajes, las colas o el estado crítico deberían replicarse o persistirse, ya que, si no hay redundancia, la torre sí es un punto único de fallo.

\subsection*{d)}

Command encapsula cada solicitud como objeto:

\begin{lstlisting}
public interface IComandoTrafico
{
    void Ejecutar();
}

public sealed class ComandoAterrizaje : IComandoTrafico
{
    private readonly Avion _avion;
    public ComandoAterrizaje(Avion avion) => _avion = avion;
    public void Ejecutar() => _avion.SolicitarAterrizaje();
}

public sealed class ComandoDespegue : IComandoTrafico
{
    private readonly Avion _avion;
    public ComandoDespegue(Avion avion) => _avion = avion;
    public void Ejecutar() => _avion.SolicitarDespegue();
}

public sealed class ComandoCambioAltitud : IComandoTrafico
{
    private readonly Avion _avion;
    private readonly int _altitud;
    public ComandoCambioAltitud(Avion avion, int altitud)
    {
        _avion = avion;
        _altitud = altitud;
    }
    public void Ejecutar() => _avion.CambiarAltitud(_altitud);
}
\end{lstlisting}

State modela estados del avión:

\begin{lstlisting}
public interface IEstadoAvion
{
    void SolicitarAterrizaje(Avion avion);
    void SolicitarDespegue(Avion avion);
    void CambiarAltitud(Avion avion, int altitud);
}

public sealed class EnVuelo : IEstadoAvion
{
    public void SolicitarAterrizaje(Avion avion)
    {
        avion.Enviar(TipoMensaje.Aterrizaje);
        avion.CambiarEstado(new Aterrizando());
    }

    public void SolicitarDespegue(Avion avion) =>
        Console.WriteLine($"{avion.Id} ya está en vuelo");

    public void CambiarAltitud(Avion avion, int altitud) =>
        avion.Enviar(TipoMensaje.CambioAltitud, altitud);
}

public sealed class Aterrizando : IEstadoAvion
{
    public void SolicitarAterrizaje(Avion avion) =>
        Console.WriteLine($"{avion.Id} ya está aterrizando");

    public void SolicitarDespegue(Avion avion) =>
        Console.WriteLine($"{avion.Id} no puede despegar mientras aterriza");

    public void CambiarAltitud(Avion avion, int altitud) =>
        Console.WriteLine($"{avion.Id} no cambia altitud mientras aterriza");
}

public sealed class EnTierra : IEstadoAvion
{
    public void SolicitarAterrizaje(Avion avion) =>
        Console.WriteLine($"{avion.Id} ya está en tierra");

    public void SolicitarDespegue(Avion avion)
    {
        avion.Enviar(TipoMensaje.Despegue);
        avion.CambiarEstado(new EnVuelo());
    }

    public void CambiarAltitud(Avion avion, int altitud) =>
        Console.WriteLine($"{avion.Id} está en tierra; no aplica cambio de altitud");
}
\end{lstlisting}

El \texttt{Avion} mantiene su estado y delega:

\begin{lstlisting}
public sealed class Avion
{
    public string Id { get; }
    private readonly IMediadorTrafico _torre;
    private IEstadoAvion _estado = new EnVuelo();

    public Avion(string id, IMediadorTrafico torre)
    {
        Id = id;
        _torre = torre;
        _torre.Registrar(this);
    }

    public void CambiarEstado(IEstadoAvion nuevo) => _estado = nuevo;
    public void Enviar(TipoMensaje tipo, int altitud = 0) =>
        _torre.Enviar(new Mensaje(Id, tipo, altitud));

    public void SolicitarAterrizaje() => _estado.SolicitarAterrizaje(this);
    public void SolicitarDespegue() => _estado.SolicitarDespegue(this);
    public void CambiarAltitud(int altitud) => _estado.CambiarAltitud(this, altitud);
}
\end{lstlisting}

Así, Command encapsula la solicitud, State decide si la solicitud es válida según el estado y Mediator centraliza la coordinación y el orden FIFO por avión.

\end{document}